% !TeX root = ../main.tex
\section{Lay Abstract}

This work examines whether modern computer vision methods for soccer analysis can be compressed to run directly on compact, low-power hardware rather than relying on costly external infrastructure. Two core tasks were explored: player and ball detection, and pitch keypoints identification. 488 model variants were trained and benchmarked, and several deployment-oriented optimisation methods were compared, including hardware acceleration, quantisation, pruning, and knowledge distillation.

Results show that hardware acceleration improved runtime substantially with almost no loss in analytical performance, while quantisation delivered the largest overall speed and energy gains but also introduced a clearer accuracy trade-off. Pruning proved more suitable for object detection than for pitch localisation, and knowledge distillation was most useful as a method for preserving accuracy in smaller models rather than as a compression method.

Finally, the best models were integrated into a full-stack pipeline that estimates possession and detects passes from broadcast video. The strongest systems operated within 10\% of other non-real-time systems, with ony 40\% GPU usage and latency under 100ms. This showa that practical on-device soccer analytics is achievable even on constrained hardware. 